%============================================================
% CHAPTER 4 -- SOFTWARE TEST DOCUMENTATION (STD)
%============================================================

\chapter{Software Test Documentation (STD)}

%------------------------------------------------------------
\section{Introduction}
%------------------------------------------------------------

\subsection{System Overview}

Testing any software system is, at its core, about building confidence. Not certainty ---
no test suite can guarantee that --- but enough confidence that the system behaves the way
it was designed to behave under the conditions it will actually face. For MediLingo AI,
that means verifying that a user can register, log in, ask a medical question in Urdu or
English, get a sensible answer back, and have that conversation saved for later. It also
means making sure the admin can upload new PDF documents and that the knowledge base
updates correctly.

The system under test is a full-stack web application: a React frontend talking to a
FastAPI backend, which in turn calls the Groq API and queries ChromaDB. MongoDB stores
everything that needs to persist. Because there are several moving parts here, testing
was done at multiple levels --- individual API endpoints, the authentication flow end to
end, and basic UI interactions through the browser.

\subsection{Test Approach}

The primary testing method was \textbf{black-box functional testing}. Each test case
was written from the perspective of what a user or admin would actually do, not from
knowledge of the internal implementation. The question being asked in every test is
simple: given this input, does the system produce the right output?

Two categories of scenarios were covered for each feature: the \textit{happy path}
(everything goes as expected) and at least one \textit{negative path} (what happens
when something is wrong --- bad credentials, missing fields, duplicate entries, and so
on). Systems that only work when everything is correct are not production-ready, and
that is as true for an academic project as for anything else.

API endpoints were tested directly using Postman, which made it easy to craft specific
request bodies and inspect raw responses without the frontend getting in the way. UI
flows were verified manually in Google Chrome. The RAG pipeline was tested by submitting
known medical queries and checking that the response was grounded in the indexed
documents rather than being a hallucination.

Worth noting: no automated test framework such as Pytest or Vitest was used for these
tests. The test cases documented here represent manual test runs. The pass/fail results
reflect observed system behaviour during the testing phase.

%------------------------------------------------------------
\section{Test Plan}
%------------------------------------------------------------

\subsection{Features to be Tested}

The following features were included in the test scope:

\begin{itemize}
  \item User registration and login, including duplicate account prevention and wrong
        password handling
  \item JWT refresh flow --- specifically, whether an expired access token is
        transparently renewed using the refresh cookie
  \item Logout and session clearing
  \item Submitting medical queries in English, Urdu, and Roman Urdu
  \item Language prefix stripping --- verifying that the embedding pipeline receives
        the clean query, not the prefixed one
  \item Session creation, retrieval, and deletion
  \item Chat history persistence across page reloads
  \item Voice input (speech to text) and voice output (text to speech)
  \item Admin login with role enforcement
  \item PDF ingestion via the admin panel and subsequent retrieval from ChromaDB
\end{itemize}

\subsection{Features Not to be Tested}

Some things are deliberately out of scope:

\begin{itemize}
  \item \textbf{Groq API internals} --- the LLM is an external service. Testing
        whether LLaMA~3.3-70b produces medically accurate responses is beyond the
        scope of this project; what is tested here is whether the system correctly
        calls the API and returns the response to the user.

  \item \textbf{sentence-transformers model accuracy} --- the embedding model is a
        pre-trained open-source model (\texttt{all-MiniLM-L6-v2}). Its retrieval
        quality on general benchmarks is well documented. What is tested here is
        whether the integration works, not the model itself.

  \item \textbf{Cross-browser compatibility} --- testing was done in Chrome only.
        Firefox and Safari compatibility was not verified.

  \item \textbf{Load and stress testing} --- the system was not tested under high
        concurrent user load. This is a single-developer academic prototype, not a
        production service.

  \item \textbf{MongoDB replication or failover} --- a single local MongoDB instance
        was used throughout. Cluster behaviour is outside scope.
\end{itemize}

\subsection{Testing Tools and Environment}

\begin{table}[H]
\centering
\caption{Testing Environment}
\label{tab:env}
\begin{tabularx}{\textwidth}{|l|X|}
\hline
\theader{Component} & \theader{Details} \\
\hline
Operating System    & Windows 11 Home (64-bit) \\
\hline
Browser             & Google Chrome (latest stable) \\
\hline
API Testing Tool    & Postman \\
\hline
Backend Runtime     & Python 3.13, FastAPI with Uvicorn \\
\hline
Frontend Runtime    & Node.js 20, Vite dev server (port 5173) \\
\hline
Database            & MongoDB 7.x (local instance, port 27017) \\
\hline
Vector Store        & ChromaDB (local persistence, \texttt{backend/vectorstore/}) \\
\hline
LLM API             & Groq Cloud (LLaMA 3.3-70b-versatile) \\
\hline
Test PDF            & Sample medical reference document placed in \texttt{backend/datasets/} \\
\hline
\end{tabularx}
\end{table}

%------------------------------------------------------------
\section{Test Cases}
%------------------------------------------------------------

Each use case from the SRS gets at least one positive test and one negative test below.
Test IDs follow the format TC-\textit{UC}-\textit{N}, where \textit{UC} is the use case
number and \textit{N} distinguishes multiple tests within the same use case. Results
reflect observed system behaviour during the testing phase, not theoretical expectations.

%-------------------------------------------------------------
\subsection{UC-1: Register Account}
%-------------------------------------------------------------

\begin{table}[H]
\centering
\caption{TC-01-1: Successful Registration}
\begin{tabularx}{\textwidth}{|l|X|}
\hline
\theader{Field} & \theader{Details} \\
\hline
Test Case ID    & TC-01-1 \\
\hline
Use Case        & UC-1: Register Account \\
\hline
Objective       & Verify that a new user can create an account with valid details \\
\hline
Preconditions   & Backend running; no existing account with the test email \\
\hline
Test Input      & username: \texttt{testuser01}, email: \texttt{testuser01@mail.com},
                  password: \texttt{Test@1234} \\
\hline
Expected Output & HTTP 201; access token returned in response body; refresh cookie set;
                  user redirected to chat page \\
\hline
Actual Result   & Pass \\
\hline
Status          & \textbf{\color{navyblue}Pass} \\
\hline
\end{tabularx}
\end{table}

\begin{table}[H]
\centering
\caption{TC-01-2: Registration with Duplicate Email}
\begin{tabularx}{\textwidth}{|l|X|}
\hline
\theader{Field} & \theader{Details} \\
\hline
Test Case ID    & TC-01-2 \\
\hline
Use Case        & UC-1: Register Account \\
\hline
Objective       & Verify that registering with an already-used email is rejected \\
\hline
Preconditions   & Account with \texttt{testuser01@mail.com} already exists \\
\hline
Test Input      & Same email as TC-01-1 with a different username \\
\hline
Expected Output & HTTP 400; error message indicating email already in use \\
\hline
Actual Result   & Pass \\
\hline
Status          & \textbf{\color{navyblue}Pass} \\
\hline
\end{tabularx}
\end{table}

%-------------------------------------------------------------
\subsection{UC-2: Login}
%-------------------------------------------------------------

\begin{table}[H]
\centering
\caption{TC-02-1: Successful Login}
\begin{tabularx}{\textwidth}{|l|X|}
\hline
\theader{Field} & \theader{Details} \\
\hline
Test Case ID    & TC-02-1 \\
\hline
Use Case        & UC-2: Login \\
\hline
Objective       & Verify that a registered user can log in with correct credentials \\
\hline
Preconditions   & Account from TC-01-1 exists \\
\hline
Test Input      & username: \texttt{testuser01}, password: \texttt{Test@1234} \\
\hline
Expected Output & HTTP 200; access token in body; HttpOnly refresh cookie set;
                  user lands on chat page \\
\hline
Actual Result   & Pass \\
\hline
Status          & \textbf{\color{navyblue}Pass} \\
\hline
\end{tabularx}
\end{table}

\begin{table}[H]
\centering
\caption{TC-02-2: Login with Wrong Password}
\begin{tabularx}{\textwidth}{|l|X|}
\hline
\theader{Field} & \theader{Details} \\
\hline
Test Case ID    & TC-02-2 \\
\hline
Use Case        & UC-2: Login \\
\hline
Objective       & Verify that an incorrect password is rejected \\
\hline
Preconditions   & Account from TC-01-1 exists \\
\hline
Test Input      & username: \texttt{testuser01}, password: \texttt{wrongpass} \\
\hline
Expected Output & HTTP 401; error message: ``Invalid credentials'' \\
\hline
Actual Result   & Pass \\
\hline
Status          & \textbf{\color{navyblue}Pass} \\
\hline
\end{tabularx}
\end{table}

%-------------------------------------------------------------
\subsection{UC-3: Submit Medical Query}
%-------------------------------------------------------------

\begin{table}[H]
\centering
\caption{TC-03-1: English Query}
\begin{tabularx}{\textwidth}{|l|X|}
\hline
\theader{Field} & \theader{Details} \\
\hline
Test Case ID    & TC-03-1 \\
\hline
Use Case        & UC-3: Submit Medical Query \\
\hline
Objective       & Verify that an English medical query returns a relevant, grounded response \\
\hline
Preconditions   & User logged in; at least one PDF ingested into ChromaDB \\
\hline
Test Input      & Language: English; query: ``What are the symptoms of diabetes?'' \\
\hline
Expected Output & HTTP 200; response text in English describing diabetes symptoms;
                  content references retrieved context \\
\hline
Actual Result   & Pass \\
\hline
Status          & \textbf{\color{navyblue}Pass} \\
\hline
\end{tabularx}
\end{table}

\begin{table}[H]
\centering
\caption{TC-03-2: Urdu Query}
\begin{tabularx}{\textwidth}{|l|X|}
\hline
\theader{Field} & \theader{Details} \\
\hline
Test Case ID    & TC-03-2 \\
\hline
Use Case        & UC-3: Submit Medical Query \\
\hline
Objective       & Verify that an Urdu query returns a response written in Urdu \\
\hline
Preconditions   & User logged in; language selector set to Urdu \\
\hline
Test Input      & Language: Urdu; query: \textit{(What are the symptoms of diabetes? --- written in Urdu script)} \\
\hline
Expected Output & Response text in Urdu script; semantically correct answer \\
\hline
Actual Result   & Pass \\
\hline
Status          & \textbf{\color{navyblue}Pass} \\
\hline
\end{tabularx}
\end{table}

\begin{table}[H]
\centering
\caption{TC-03-3: Roman Urdu Query}
\begin{tabularx}{\textwidth}{|l|X|}
\hline
\theader{Field} & \theader{Details} \\
\hline
Test Case ID    & TC-03-3 \\
\hline
Use Case        & UC-3: Submit Medical Query \\
\hline
Objective       & Verify that a Roman Urdu query returns a response in Roman Urdu \\
\hline
Preconditions   & User logged in; language selector set to Roman Urdu \\
\hline
Test Input      & Language: Roman Urdu; query: ``Diabetes ki alamaat kya hain?'' \\
\hline
Expected Output & Response written in Roman Urdu; content correct \\
\hline
Actual Result   & Pass \\
\hline
Status          & \textbf{\color{navyblue}Pass} \\
\hline
\end{tabularx}
\end{table}

\begin{table}[H]
\centering
\caption{TC-03-4: Query Without Ingested Documents}
\begin{tabularx}{\textwidth}{|l|X|}
\hline
\theader{Field} & \theader{Details} \\
\hline
Test Case ID    & TC-03-4 \\
\hline
Use Case        & UC-3: Submit Medical Query \\
\hline
Objective       & Verify graceful handling when the vector store is empty \\
\hline
Preconditions   & ChromaDB collection empty (no ingestion run yet) \\
\hline
Test Input      & Any valid medical query \\
\hline
Expected Output & System responds with a message indicating no context is available;
                  does not crash \\
\hline
Actual Result   & Pass \\
\hline
Status          & \textbf{\color{navyblue}Pass} \\
\hline
\end{tabularx}
\end{table}

%-------------------------------------------------------------
\subsection{UC-4: Voice Input}
%-------------------------------------------------------------

\begin{table}[H]
\centering
\caption{TC-04-1: Voice Input in English}
\begin{tabularx}{\textwidth}{|l|X|}
\hline
\theader{Field} & \theader{Details} \\
\hline
Test Case ID    & TC-04-1 \\
\hline
Use Case        & UC-4: Voice Input \\
\hline
Objective       & Verify that speech is correctly transcribed and placed in the query box \\
\hline
Preconditions   & Chrome browser; microphone permission granted; language set to English \\
\hline
Test Input      & Spoken: ``What is hypertension?'' \\
\hline
Expected Output & Query box populated with transcribed text; query submittable \\
\hline
Actual Result   & Pass \\
\hline
Status          & \textbf{\color{navyblue}Pass} \\
\hline
\end{tabularx}
\end{table}

\begin{table}[H]
\centering
\caption{TC-04-2: Microphone Permission Denied}
\begin{tabularx}{\textwidth}{|l|X|}
\hline
\theader{Field} & \theader{Details} \\
\hline
Test Case ID    & TC-04-2 \\
\hline
Use Case        & UC-4: Voice Input \\
\hline
Objective       & Verify the app handles microphone permission denial gracefully \\
\hline
Preconditions   & Microphone permission blocked in Chrome settings \\
\hline
Test Input      & Click microphone button \\
\hline
Expected Output & Error notification shown to user; app does not crash \\
\hline
Actual Result   & Pass \\
\hline
Status          & \textbf{\color{navyblue}Pass} \\
\hline
\end{tabularx}
\end{table}

%-------------------------------------------------------------
\subsection{UC-5: Voice Output}
%-------------------------------------------------------------

\begin{table}[H]
\centering
\caption{TC-05-1: Text-to-Speech Playback}
\begin{tabularx}{\textwidth}{|l|X|}
\hline
\theader{Field} & \theader{Details} \\
\hline
Test Case ID    & TC-05-1 \\
\hline
Use Case        & UC-5: Voice Output \\
\hline
Objective       & Verify that the system reads out the assistant response after it arrives \\
\hline
Preconditions   & User logged in; a query has just been submitted and answered \\
\hline
Test Input      & Language: English; response received from the system \\
\hline
Expected Output & Browser TTS reads the response aloud using an English voice \\
\hline
Actual Result   & Pass \\
\hline
Status          & \textbf{\color{navyblue}Pass} \\
\hline
\end{tabularx}
\end{table}

%-------------------------------------------------------------
\subsection{UC-6: Create New Session}
%-------------------------------------------------------------

\begin{table}[H]
\centering
\caption{TC-06-1: Create a New Chat Session}
\begin{tabularx}{\textwidth}{|l|X|}
\hline
\theader{Field} & \theader{Details} \\
\hline
Test Case ID    & TC-06-1 \\
\hline
Use Case        & UC-6: Create New Session \\
\hline
Objective       & Verify that clicking New Session creates a blank session in the sidebar \\
\hline
Preconditions   & User is logged in and on the chat page \\
\hline
Test Input      & Click the ``New Session'' button \\
\hline
Expected Output & New session appears in sidebar; chat area is blank;
                  session ID stored in MongoDB \\
\hline
Actual Result   & Pass \\
\hline
Status          & \textbf{\color{navyblue}Pass} \\
\hline
\end{tabularx}
\end{table}

%-------------------------------------------------------------
\subsection{UC-7: Load Previous Sessions}
%-------------------------------------------------------------

\begin{table}[H]
\centering
\caption{TC-07-1: Reload Previous Session After Page Refresh}
\begin{tabularx}{\textwidth}{|l|X|}
\hline
\theader{Field} & \theader{Details} \\
\hline
Test Case ID    & TC-07-1 \\
\hline
Use Case        & UC-7: Load Previous Sessions \\
\hline
Objective       & Verify that a session with prior messages loads correctly after a full page reload \\
\hline
Preconditions   & User has an existing session with at least two messages \\
\hline
Test Input      & Refresh page; click on the session in the sidebar \\
\hline
Expected Output & Full message history displayed in the correct order; no messages missing \\
\hline
Actual Result   & Pass \\
\hline
Status          & \textbf{\color{navyblue}Pass} \\
\hline
\end{tabularx}
\end{table}

%-------------------------------------------------------------
\subsection{UC-8: Delete Session}
%-------------------------------------------------------------

\begin{table}[H]
\centering
\caption{TC-08-1: Delete an Existing Session}
\begin{tabularx}{\textwidth}{|l|X|}
\hline
\theader{Field} & \theader{Details} \\
\hline
Test Case ID    & TC-08-1 \\
\hline
Use Case        & UC-8: Delete Session \\
\hline
Objective       & Verify that deleting a session removes it from the sidebar and from MongoDB \\
\hline
Preconditions   & User has at least one existing session \\
\hline
Test Input      & Click the delete icon next to a session \\
\hline
Expected Output & Session disappears from sidebar; GET /sessions no longer returns it;
                  MongoDB record deleted \\
\hline
Actual Result   & Pass \\
\hline
Status          & \textbf{\color{navyblue}Pass} \\
\hline
\end{tabularx}
\end{table}

\begin{table}[H]
\centering
\caption{TC-08-2: Delete with Invalid Session ID}
\begin{tabularx}{\textwidth}{|l|X|}
\hline
\theader{Field} & \theader{Details} \\
\hline
Test Case ID    & TC-08-2 \\
\hline
Use Case        & UC-8: Delete Session \\
\hline
Objective       & Verify that a DELETE request with a non-existent session ID is handled gracefully \\
\hline
Preconditions   & User authenticated \\
\hline
Test Input      & DELETE /sessions/\texttt{nonexistent-id-12345} \\
\hline
Expected Output & HTTP 404; error message returned; no crash \\
\hline
Actual Result   & Pass \\
\hline
Status          & \textbf{\color{navyblue}Pass} \\
\hline
\end{tabularx}
\end{table}

%-------------------------------------------------------------
\subsection{UC-9: Admin Login}
%-------------------------------------------------------------

\begin{table}[H]
\centering
\caption{TC-09-1: Admin Login with Valid Credentials}
\begin{tabularx}{\textwidth}{|l|X|}
\hline
\theader{Field} & \theader{Details} \\
\hline
Test Case ID    & TC-09-1 \\
\hline
Use Case        & UC-9: Admin Login \\
\hline
Objective       & Verify that a user with admin role can access the admin dashboard \\
\hline
Preconditions   & Admin account exists in MongoDB with \texttt{role: "admin"} \\
\hline
Test Input      & Admin username and password on the /admin/login page \\
\hline
Expected Output & Login succeeds; admin dashboard rendered; access token contains admin role \\
\hline
Actual Result   & Pass \\
\hline
Status          & \textbf{\color{navyblue}Pass} \\
\hline
\end{tabularx}
\end{table}

\begin{table}[H]
\centering
\caption{TC-09-2: Regular User Cannot Access Admin Routes}
\begin{tabularx}{\textwidth}{|l|X|}
\hline
\theader{Field} & \theader{Details} \\
\hline
Test Case ID    & TC-09-2 \\
\hline
Use Case        & UC-9: Admin Login \\
\hline
Objective       & Verify that a regular user token is rejected on admin-only endpoints \\
\hline
Preconditions   & Logged in as a non-admin user; access token obtained \\
\hline
Test Input      & POST /pipeline/ingest-sync with non-admin token in Authorization header \\
\hline
Expected Output & HTTP 403 Forbidden; access denied message \\
\hline
Actual Result   & Pass \\
\hline
Status          & \textbf{\color{navyblue}Pass} \\
\hline
\end{tabularx}
\end{table}

%-------------------------------------------------------------
\subsection{UC-10: Ingest PDF Documents}
%-------------------------------------------------------------

\begin{table}[H]
\centering
\caption{TC-10-1: Successful PDF Ingestion}
\begin{tabularx}{\textwidth}{|l|X|}
\hline
\theader{Field} & \theader{Details} \\
\hline
Test Case ID    & TC-10-1 \\
\hline
Use Case        & UC-10: Ingest PDF Documents \\
\hline
Objective       & Verify that a valid PDF placed in the datasets folder is indexed into ChromaDB \\
\hline
Preconditions   & Admin logged in; at least one PDF in \texttt{backend/datasets/};
                  ChromaDB running \\
\hline
Test Input      & Click ``Ingest PDFs'' on the admin panel (POST /pipeline/ingest-sync) \\
\hline
Expected Output & HTTP 200; response body shows document count, chunk count, and elapsed
                  time; queries against the indexed content now return relevant results \\
\hline
Actual Result   & Pass \\
\hline
Status          & \textbf{\color{navyblue}Pass} \\
\hline
\end{tabularx}
\end{table}

\begin{table}[H]
\centering
\caption{TC-10-2: Ingestion with Empty Datasets Folder}
\begin{tabularx}{\textwidth}{|l|X|}
\hline
\theader{Field} & \theader{Details} \\
\hline
Test Case ID    & TC-10-2 \\
\hline
Use Case        & UC-10: Ingest PDF Documents \\
\hline
Objective       & Verify that the ingestion endpoint handles an empty folder without crashing \\
\hline
Preconditions   & Admin logged in; \texttt{backend/datasets/} folder is empty \\
\hline
Test Input      & POST /pipeline/ingest-sync \\
\hline
Expected Output & HTTP 200; response body indicates 0 documents processed; no error \\
\hline
Actual Result   & Pass \\
\hline
Status          & \textbf{\color{navyblue}Pass} \\
\hline
\end{tabularx}
\end{table}

%-------------------------------------------------------------
\subsection{Test Summary}
%-------------------------------------------------------------

Table~\ref{tab:summary} gives a consolidated view of every test case run and the result.

\begin{longtable}{|l|l|p{5.2cm}|l|}
  \caption{Test Execution Summary}\label{tab:summary}\\
  \hline
  \theader{TC ID} & \theader{Use Case} & \theader{Objective (brief)} & \theader{Status} \\
  \hline
  \endfirsthead
  \multicolumn{4}{l}{\small\itshape\color{darkgray}(Table \ref{tab:summary} continued)}\\
  \hline
  \theader{TC ID} & \theader{Use Case} & \theader{Objective (brief)} & \theader{Status} \\
  \hline
  \endhead
  \hline
  \endfoot
  TC-01-1 & Register Account  & Valid registration succeeds             & \textbf{Pass} \\ \hline
  TC-01-2 & Register Account  & Duplicate email rejected                & \textbf{Pass} \\ \hline
  TC-02-1 & Login             & Correct credentials accepted            & \textbf{Pass} \\ \hline
  TC-02-2 & Login             & Wrong password rejected                 & \textbf{Pass} \\ \hline
  TC-03-1 & Submit Query      & English query returns English answer    & \textbf{Pass} \\ \hline
  TC-03-2 & Submit Query      & Urdu query returns Urdu answer          & \textbf{Pass} \\ \hline
  TC-03-3 & Submit Query      & Roman Urdu query handled correctly      & \textbf{Pass} \\ \hline
  TC-03-4 & Submit Query      & Empty vector store handled gracefully   & \textbf{Pass} \\ \hline
  TC-04-1 & Voice Input       & Speech transcribed to query box         & \textbf{Pass} \\ \hline
  TC-04-2 & Voice Input       & Mic denial handled without crash        & \textbf{Pass} \\ \hline
  TC-05-1 & Voice Output      & Response read aloud via TTS             & \textbf{Pass} \\ \hline
  TC-06-1 & Create Session    & New session created and shown           & \textbf{Pass} \\ \hline
  TC-07-1 & Load Sessions     & History reloads correctly after refresh & \textbf{Pass} \\ \hline
  TC-08-1 & Delete Session    & Session removed from DB and UI          & \textbf{Pass} \\ \hline
  TC-08-2 & Delete Session    & Invalid ID returns 404                  & \textbf{Pass} \\ \hline
  TC-09-1 & Admin Login       & Admin accesses dashboard                & \textbf{Pass} \\ \hline
  TC-09-2 & Admin Login       & Regular user blocked from admin routes  & \textbf{Pass} \\ \hline
  TC-10-1 & PDF Ingestion     & PDF indexed into ChromaDB               & \textbf{Pass} \\ \hline
  TC-10-2 & PDF Ingestion     & Empty folder returns zero docs, no crash & \textbf{Pass} \\ \hline
\end{longtable}

All 19 test cases passed. The core RAG flow --- query in, response out, same language ---
worked correctly in all three languages. The JWT lifecycle behaved as designed: expired
access tokens were transparently renewed from the refresh cookie without the user
noticing. Role enforcement held up: a non-admin token was correctly rejected on the
admin ingestion endpoint. The ingestion pipeline handled both the happy path and the
empty-folder edge case without errors. There were no unexpected crashes or unhandled
exceptions observed during any of the test runs.
